% ============================================================
%  Chapter 4 — Competitive Analysis
% ============================================================
\chapter{Analysis of Similar Components}
\label{ch:competitive-analysis}

This chapter surveys commercially and open-source components that solve the
same or overlapping problems as \textbf{resilient-call}, examines how they are
presented and documented, and uses these findings to evaluate and improve
this component.

\section{Components Surveyed}
\label{sec:components-surveyed}

\begin{table}[H]
    \centering
    \caption{Overview of surveyed components}
    \label{tab:competitors}
    \begin{tabularx}{\linewidth}{>{\ttfamily}l l l X}
        \toprule
        \textbf{Name} & \textbf{Language} & \textbf{License} & \textbf{Features} \\
        \midrule
        tenacity      & Python & Apache 2.0 & Retry only; decorator + functional API \\
        pybreaker      & Python & BSD         & Circuit breaker only; basic API \\
        Resilience4j  & Java   & Apache 2.0 & Retry + CB + bulkhead + rate-limiter + time-limiter \\
        Polly         & C\#/.NET & BSD       & Retry + CB + timeout + bulkhead + cache + fallback \\
        \bottomrule
    \end{tabularx}
\end{table}

\section{Commercial and Open-Source Presentation}
\label{sec:presentation-research}

\subsection{tenacity (Python)}

\textbf{tenacity}~\cite{tenacity} is the most widely used Python retry library
(over 20 million monthly PyPI downloads).  Its strengths from a documentation
perspective:

\begin{itemize}[leftmargin=*, label=\textbullet]
    \item Concise README with runnable code examples grouped by use case.
    \item Decorator, context manager, and functional styles all demonstrated.
    \item \texttt{retry\_if\_exception\_type}, \texttt{retry\_if\_result},
          and \texttt{wait\_*} composables provide powerful declarative
          configuration without subclassing.
    \item \textbf{No circuit breaker} — the library explicitly scopes itself
          to retry only.
\end{itemize}

\subsection{pybreaker (Python)}

\textbf{pybreaker}~\cite{pybreaker} provides circuit breaker functionality
but deliberately omits retry.  Notable characteristics:

\begin{itemize}[leftmargin=*, label=\textbullet]
    \item Minimal API — only one class (\texttt{CircuitBreaker}) and one
          decorator.
    \item Documentation is limited to a single README with one code example.
    \item Thread-safe, but no async support.
    \item Supports state listeners (equivalent to \texttt{on\_state\_change}).
    \item \textbf{No retry} and \textbf{no fallback} mechanism.
\end{itemize}

\subsection{Resilience4j (Java)}

\textbf{Resilience4j}~\cite{resilience4j} is the de facto standard in the
Java ecosystem, inspired by Netflix Hystrix.  It sets the documentation
gold standard:

\begin{itemize}[leftmargin=*, label=\textbullet]
    \item Dedicated documentation website with separate pages per module.
    \item Architecture diagrams, state machine diagrams, and event-flow
          diagrams embedded in the documentation.
    \item Extensive configuration tables with types, defaults, and constraints.
    \item Javadoc generated and published alongside the user guide.
    \item \textbf{Spring Boot autoconfiguration} and metrics integration
          (Micrometer, Prometheus) documented separately.
    \item Versioned release notes.
\end{itemize}

\subsection{Polly (.NET)}

\textbf{Polly}~\cite{polly} provides the richest feature set of all surveyed
libraries.  Documentation highlights:

\begin{itemize}[leftmargin=*, label=\textbullet]
    \item GitHub wiki with one page per policy type, linked from the README.
    \item Policy composition patterns documented with code and narrative.
    \item Telemetry and resilience events documented, with integration guides
          for Microsoft.Extensions.Http.
    \item Migration guides between major versions.
    \item Community-maintained FAQ and ``anti-patterns'' page.
\end{itemize}

\section{Documentation Comparison}
\label{sec:documentation-comparison}

\begin{table}[H]
    \centering
    \caption{Documentation feature comparison}
    \label{tab:doc-comparison}
    \begin{tabularx}{\linewidth}{X >{\centering\arraybackslash}p{1.8cm} >{\centering\arraybackslash}p{1.8cm} >{\centering\arraybackslash}p{1.8cm} >{\centering\arraybackslash}p{1.8cm} >{\centering\arraybackslash}p{2.2cm}}
        \toprule
        \textbf{Feature} & \textbf{tenacity} & \textbf{pybreaker} & \textbf{Resilience4j} & \textbf{Polly} & \textbf{resilient-call} \\
        \midrule
        API reference (generated) & \checkmark & -- & \checkmark & \checkmark & \checkmark \\
        Quickstart ($\leq$ 5 min)  & \checkmark & \checkmark & \checkmark & \checkmark & \checkmark \\
        Configuration table         & --         & --          & \checkmark & \checkmark & \checkmark \\
        State / architecture diagrams & --       & --          & \checkmark & --         & \checkmark \\
        Code examples ($\geq$ 3)   & \checkmark & --          & \checkmark & \checkmark & \checkmark \\
        Changelog                  & \checkmark & \checkmark  & \checkmark & \checkmark & \checkmark \\
        License stated             & \checkmark & \checkmark  & \checkmark & \checkmark & \checkmark \\
        Async documented           & \checkmark & --          & \checkmark & \checkmark & \checkmark \\
        Fallback / degradation     & --         & --          & \checkmark & \checkmark & \checkmark \\
        Anti-patterns / warnings   & --         & --          & --         & \checkmark & \checkmark \\
        \bottomrule
    \end{tabularx}
\end{table}

\section{Findings and Improvements Applied}
\label{sec:findings}

Based on the research above, the following improvements were made to
\textbf{resilient-call}:

\begin{enumerate}[leftmargin=*, label=\arabic*.]
    \item \textbf{Configuration tables} (inspired by Resilience4j and Polly).
          Every class's parameters are documented in a table with type, default,
          and description — not just prose.

    \item \textbf{Architecture and state diagrams} (inspired by Resilience4j).
          pybreaker and tenacity include no diagrams; Resilience4j's
          state-machine diagram is one of its most cited pages.  Diagrams were
          produced for the class structure, state machine, and three sequence
          scenarios.

    \item \textbf{Idempotency warning} (inspired by Polly's anti-patterns
          page).  A prominently displayed warning was added to the user guide
          explaining the conditions under which retry is unsafe.

    \item \textbf{5-minute quickstart example} (matching tenacity's
          approachable README style).  \texttt{examples/quickstart.py} is a
          self-contained runnable script demonstrating all three usage
          patterns.

    \item \textbf{Combined retry + circuit breaker} (gap identified in both
          tenacity and pybreaker).  Neither library offers both mechanisms;
          users must compose them manually with adapter code.
          \texttt{ResilientExecutor} solves this directly.

    \item \textbf{Deployment restrictions documented explicitly} (inspired by
          Polly's multi-host documentation).  The limitation that circuit
          breaker state is not shared across processes is called out in both
          the architecture chapter and the user guide.
\end{enumerate}

\section{Summary}
\label{sec:competitive-summary}

\textbf{resilient-call} occupies a niche that is currently underserved in the
Python ecosystem: a single library that provides both retry and circuit-breaker
functionality in a composable, decorator-friendly API.  The closest equivalent
in the Java ecosystem is Resilience4j, whose documentation and architecture
transparency served as the primary benchmark for this component's own
documentation quality.
